% 本章原书的节号是阿拉伯数字（1, 2, 3, …），与全书默认的罗马数字不同。
% 按 STYLE_GUIDE §3.2「各章按原书自己的节号体例」：章前局部改写，章末 reflist 之后恢复。
\renewcommand{\thesection}{\arabic{section}}
% ===== PDF p.068 (原书 p.60) =====
\chapter{奇点与拓扑场论}
\chapterauthor{高橋篤史}{[大阪大学大学院理学研究科]}

\section{引言}

超弦理论是一种物理理论：它认为基本粒子与空间都由具有 1 维自由度的弦的振动、运动来描述；作为“万物理论”的有力候选，它正被人们积极研究。它与数学的各个领域也密切相关，给数学与物理双方都带来了巨大影响。其中，主张“两种不同的几何学总是成对存在”的镜像对称，正是数学与超弦理论关联最强的地方。

标题里的另一个关键词“奇点”，各位会联想到什么样的几何对象呢？查一查“奇点”这个条目，词典上写着\footnote{我很好奇能否不用数学定义、只用平易的日常用语来写，于是查了词典。}：“曲线或曲面上，在该点处切线或切平面不存在、或者有两条以上的切线或切平面的点”。这话确实不假，不过这里请把它理解为更接近“特异”的一个意思：“与周围的结构相比性质格外不同、具有特别优异性质的点”。只要限定在某一类好的奇点上来考虑，想必各位马上就能认同：这种特殊的点凝聚着多么丰富的对称性与整体结构的信息。镜像对称中出现的正是这种特殊点；例如，本来光滑的非交换流形，正因为我们是透过交换几何（复代数簇）去看它，它的性质才以奇异性的形式显现出来——也可以这样解释。% [Tutor Note: 此句是说：非交换几何中的空间本来“光滑”，一旦用交换几何（复代数簇）来观察，其非交换性便表现为奇点。这里“非交换流形”是非交换几何中的空间，并非字面上的“非交换的流形”。]

本书将在本章和第 11 章解说镜像对称与奇点的关系。这里考虑的奇点，是出现在由一个多项式的零点集合所表示的图形上的特殊点，即超曲面奇点。本章讲解与经典镜像对称、也就是拓扑弦理论的镜像对称相关的话题，第 11 章则讲解与同调镜像对称相关的话题。

% ===== PDF p.069 (原书 p.61) =====
\section{超弦理论与奇点的相遇}

首先，我想简要回顾一下奇点理论与超弦理论之间关系的历史。关于超弦理论的更详细解说，请参阅第 1 章。

众所周知，超弦理论只有在时空为 10 维时才能自洽地定义。另一方面，我们居住的时空是 4 维的。因此有一种被称为“超弦理论的紧化”的想法：假定这个 10 维时空分为 4 维与 6 维的空间，6 维空间非常小、看不见，所以平时只认识到剩下的 4 维时空，并在此假定下构建理论。按照这一想法，我们要得到支配 4 维时空中电子、光子、夸克等基本粒子的统一理论，即 N=1 超对称场论。为此，需要让弦的世界面以这个 6 维空间为目标空间，并在世界面上考虑具有 N=2 超对称性的共形场论。而且，为了让这个 N=2 共形场论能够自洽地定义，就要求 6 维空间是 Calabi--Yau 流形。

这样，N=2 共形场论的研究就变得重要起来。然而，像上面那样从超弦理论经由 Calabi--Yau 流形构造出的 N=2 共形场论（称为 sigma 模型）非常困难。它并不太适合用来具体计算各种量。因此需要另想办法来构造和研究理论。于是，人们引入了以超曲面奇点为出发点的 N=2 共形场论，也就是 N=2 Landau--Ginzburg 理论。这个理论只从被称为超势的一个多项式\footnote{一般来说是全纯函数。}所携带的信息出发，试图得到理论的真空、状态、相互作用等全部物理量。

从 80 年代后半期起，人们对 N=2 Landau--Ginzburg 理论作了热烈研究。但在数学一侧，同样的研究对象出于完全不同的动机，早在 1980 年前后就已经出现了。这就是斋藤恭司基于原始形式的周期映射理论。

% ===== PDF p.070 (原书 p.62) =====
斋藤的目标是把椭圆积分论一般化、高维化。他把椭圆曲线的族加以推广，考察超曲面奇点及其半通用形变，并把第一类微分的自然推广——原始形式——用若干条公理定义出来。然后，他证明了在加权齐次超曲面奇点的情况下原始形式存在\footnote{一般情形由斋藤盛彦证明。}，且半通用形变的底空间（掌管形变的那些参数构成的空间）上具有平坦结构。

把这里叙述的斋藤的结果用 N=2 Landau--Ginzburg 理论的语言改写，就是下面这段话。对给定的超势施加微扰，考察物理系统因该微扰而如何变化；由此证明理论中存在真空\footnote{物理学中真空的存在总是被假定；但这本来是必须证明的主张。}，并计算由微扰空间上确定的 0 亏格关联函数。

从下一节起，我们用公式更具体地看一看。首先，我们来看一看从多项式出发最终想得到的东西——平坦结构与关联函数——究竟是什么样子。做法和第 2 章一样，以复射影空间的量子上同调为例来说明。

\section{量子上同调环的复习}

首先，设 $\CP^n$ 为 $n$ 维复射影空间。它是由 $n+1$ 个复数之比 $(x_0:\cdots:x_n)$ 构成的空间\footnote{也可以用 $n+1$ 张 $\C^n$ 粘贴起来作出。}。因此，$\CP^{n-1}$ 可以写成 $(x_0:\cdots:x_{n-1}:0)$，成为 $\CP^n$ 的 $n-1$ 维子空间；依此类推，就得到子空间的链 $\CP^n\supset\CP^{n-1}\supset\cdots\supset\CP^1\supset\CP^0=(1:0)$。粗略地说，对 $\CP^n$ 中的这些子空间，把能通过连续形变互相移到的看作同一个。这样作出的阿贝尔群就是同调群 $H_*(\CP^n)$，它的 Poincaré 对偶就是上同调群 $H^*(\CP^n)$。在本例中，用上面作出的子空间，把 $H^*(\CP^n)$ 按直和分解 $H^*(\CP^n)\cong\bigoplus_{i=0}^{n}H^{2i}(\CP^n)$ 来看时，有
\begin{equation}
  H^{2i}(\CP^n)\cong \C[\CP^{n-i}]
  \tag{1}
\end{equation}
成立\footnote{为简单起见，系数不取 $\Z$ 而取 $\C$；并且通过 Poincaré 对偶，把上同调的元素与同调的元素等同起来书写。}。

把 $H^*(\CP^n)$ 看作 $n+1$ 维仿射空间，引入平坦坐标 $(t_0,t_1,\dots,t_n)$。特别地，让每个子空间 $H^{2i}(\CP^n)$ 分别对应 $t_i$。进而，把 $H^*(\CP^n)$ 与原点处的切空间等同起来，考虑同型
\[
  H^*(\CP^n)\cong T_0H^*(\CP^n)\cong\bigoplus_{i=0}^{n}\C\frac{\partial}{\partial t_i}.
\]

% ===== PDF p.071 (原书 p.63) =====
虽然来得有些突兀，Gromov--Witten 不变量的母函数（Gromov--Witten 势）由下式给出：
\begin{equation}
  F_0^{\CP^n}(t):=\frac{1}{6}\sum_{a+b+c=n}t_at_bt_c
  +\sum_{d>0}\sum_{k_2,\dots,k_n\ge 0}
  N_{d:k_2,\dots,k_n}\frac{t_2^{k_2}\cdots t_n^{k_n}}{k_2!\cdots k_n!}\,e^{dt_1}.
  \tag{2}
\end{equation}
这里 $N_{d:k_2,\dots,k_n}$ 是 $d$ 次 Gromov--Witten 不变量。也就是说，它数的是与 $\CP^n$ 中处于一般位置的 $k_i$ 个 $n-i$ 维子空间 $\CP^{n-i}$ 相交的 $d$ 次有理曲线的条数；在 $\sum_{i\ge 2}(1-i)k_i+(n+1)d=3-n$ 以外时定义为 0。请注意，$t_0$ 只出现在势函数的 3 次项中。

% [Tutor Note: 原书此处印作 $\delta_{i,n-i}$，下标只用到 $i$；但 $\eta_{ij}=C_{0ij}$ 是 $H^{2i}$ 与 $H^{2j}$ 的相交配对，非零当且仅当 $i+j=n$，故应为 $\delta_{i+j,n}$（等价于 $\delta_{j,n-i}$）。此处照录原书，读者按相交配对理解即可。]
% [Tutor Note: 原文「（コ）ホモロジー」把「同调」与「上同调」并排写出，带有玩笑意味；正文译作「（上）同调」。]
考虑 $F_0^{\CP^n}$ 的三阶微分 $C_{ijk}:=\partial F_0^{\CP^n}/\partial t_i\partial t_j\partial t_k$，由 Gromov--Witten 不变量的性质可以导出两条重要性质。第一条是 $H^*(\CP^n)$ 上的非退化双线性形式 $\eta:H^*(\CP^n)\otimes H^*(\CP^n)\to\C$，定义为 $\eta(\partial/\partial t_i,\partial/\partial t_j):=C_{0ij}=\delta_{i,n-i}$\footnote{这里 $\delta_{ab}$ 是 Kronecker delta 记号。}。在本例中，$\eta$ 恰好与（上）同调的相交数一致。另一条是 WDVV 方程
\begin{equation}
  \sum_{a,b=0}^{n}C_{ija}(\eta_{ab})^{-1}C_{bkl}
  =\sum_{a,b=0}^{n}C_{ika}(\eta_{ab})^{-1}C_{bjl}
  \tag{3}
\end{equation}
特别地，当 $n=2$ 时，这个 WDVV 方程与式
\begin{equation}
  C_{222}^{2}=C_{112}^{2}-C_{111}C_{122}
  \tag{4}
\end{equation}
等价；而 $N_{d:k_2}$ 除 $k_2=3d-1$ 以外为 0，所以令 $N_d:=N_{d:3d-1}$，就得到第 2 章中的式 (25)
\begin{equation}
  N_d=\sum_{a+b=d}N_aN_b\Bigl(a^2b^2\binom{3d-4}{3a-2}-a^3b\binom{3d-4}{3a-1}\Bigr)
  \tag{5}
\end{equation}
由于过 $\CP^2$ 中两点的直线恰好只有一条，也就是 $N_1=1$，所以所有的 $N_d$ 都能归纳地算出。

在 $H^*(\CP^n)$ 上，把由 $t$ 参数化的积\footnote{准确地说，把 $\C[[t_0,t_1,\dots,t_n]]$ 取为 $t_0,t_1,\dots,t_n$ 的形式幂级数环时，考虑的是 $H^*(\CP^n)\otimes_{\C}\C[[t_0,t_1,\dots,t_n]]$ 上的 $\C$-双线性积。}定义为
\begin{equation}
  \frac{\partial}{\partial t_i}\circ\frac{\partial}{\partial t_j}
  :=\sum_{a,b=0}^{n}C_{ija}(\eta_{ab})^{-1}\frac{\partial}{\partial t_b}
  \tag{6}
\end{equation}
WDVV 方程等价于这个积 $\circ$ 满足结合律。特别地，$\partial/\partial t_0$ 给出这个代数的单位元。此外，作为重要的关系式
\begin{equation}
  \eta\Bigl(\frac{\partial}{\partial t_i}\circ\frac{\partial}{\partial t_j},\frac{\partial}{\partial t_k}\Bigr)
  =\eta\Bigl(\frac{\partial}{\partial t_i},\frac{\partial}{\partial t_j}\circ\frac{\partial}{\partial t_k}\Bigr)
  \tag{7}
\end{equation}
成立。

又，令 $t_i=0$（$i\ne 1$），可知代数 $(H^*(\CP^n),\circ)$ 同构于 $\C[x,q]/(x^{n+1}-q)$，其中 $q:=e^{t_1}$。这样就从势函数 $F_0^{\CP^n}$ 导出了量子上同调环。

令 $f_q:=z_1+\cdots+z_n+q\cdot(z_1\cdots z_n)^{-1}$，则应注意可写成
\begin{align}
  QH^*(\CP^n)&\cong \C[x,q]/(x^{n+1}-q)
  \tag{8}\\
  &\cong \C[z_1^{\pm},\dots,z_n^{\pm},q]\Big/
  \Bigl(z_1\frac{\partial f_q}{\partial z_1},\dots,z_n\frac{\partial f_q}{\partial z_n}\Bigr)
  \tag{9}
\end{align}
这个事实在它与奇点理论的关联上很重要。

% ===== PDF p.072 (原书 p.64) =====
% [Tutor Note: 原书式（10）与 p.65 正文都印作 $t^i$；按坐标分量的写法与下文 $\deg t_i=1-i$，此处应为 $t_i$。为存原书面貌仍照录 $t^i$，读者按 $t_i$ 理解即可。]
接下来，把被称为 Euler 向量场的向量场 $E$ 定义为
\begin{equation}
  E:=\sum_{i=0}^{n}(1-i)t^i\frac{\partial}{\partial t_i}+(n+1)\frac{\partial}{\partial t_1}
  \tag{10}
\end{equation}
给出 $E$，与给坐标 $(t_0,t_1,\dots,t_n)$ 分别赋予次数
\begin{equation}
  \deg t_i=1-i,\quad i\ne 1,\qquad \deg e^{t_1}=n+1
  \tag{11}
\end{equation}
恰好对应。于是，关于这个次数，可知势函数 $F_0^{\CP^n}$ 具有次数 $3-n$，即
\begin{equation}
  EF_0^{\CP^n}=(3-n)F_0^{\CP^n}
  \tag{12}
\end{equation}
本来想给 $t_1$ 也像其他变量一样赋予次数 0。但那样一来，势函数就无法具有次数
% ===== PDF p.073 (原书 p.65) =====
$3-n$，也就是说，它不会成为具有特征值 $3-n$ 的 $E$ 的特征函数。于是我给 $E$ 加上一个奇妙的项 $(n+1)\partial/\partial t_1$ 来作调整。这件事有几何意义，下面稍作说明。

衡量流形\footnote{指把复数域上的仿射空间粘贴起来作出的图形。}（这里就是 $\CP^n$）偏离仿射空间有多远的指标，叫作 Chern 类。特别是第一 Chern 类 $c_1(\CP^n)$，它是表示流形作尺度变换时行为的量，也就是表示它在多大程度上不满足尺度不变性的量。实际上，$(n+1)\partial/\partial t_1$ 正是 $c_1(\CP^n)$ 在同型 $H^*(\CP^n)\cong\bigoplus_{i=0}^{n}\C\,\partial/\partial t_i$ 下的像。$\sum_{i=0}^{n}(1-i)t^i\partial/\partial t_i$ 从形式上看，应当能看出它是与尺度变换相关的向量场。把这个向量场与 $c_1(\CP^n)$ 的修正合在一起，即考虑 Euler 向量场 $E$，各项恰好互相抵消；于是尽管流形本身在尺度变换下行为并不好，势函数却成为关于尺度变换共变的量。% [Tutor Note: 本段开头把 Chern 类直观地说成“衡量流形偏离仿射空间有多远的指标”，原书如此；严格说，Chern 类是切丛的示性类，第一 Chern 类为零正是 Calabi--Yau 条件。]

\section{平坦结构与 Frobenius 结构}

下面把这里得到的数据整理一下。首先考虑仿射空间 $H^*(\CP^n)$ 及其切空间 $T_0H^*(\CP^n)\cong H^*(\CP^n)$，并已定义势函数 $F_0^{\CP^n}$。$F_0^{\CP^n}$ 关于 Euler 向量场 $E$ 是齐次的（次数 $3-n$）；特别是由它的三阶微分，在 $H^*(\CP^n)$ 上定下了依赖于参数的结合积结构。进而，存在特别的坐标 $t_0$；用 $F_0^{\CP^n}$ 关于它的微分，给出了 $T_0H^*(\CP^n)$ 的非退化双线性形式（相交数）$\eta$，并描述了积结构的单位元。

基本数据的四元组 $(T_0H^*(\CP^n),F_0^{\CP^n},E,\partial/\partial t_0)$，以及到此为止所述的这个四元组所满足的性质，被公理化为 Frobenius 结构或平坦结构。它是镜像对称研究中最重要的研究对象之一\footnote{名称的发现历史顺序正好相反，但非常多的时候就直接称作 Frobenius 结构。}。这正是从流形、奇点这类几何出发、最终想要得到的代数结构。特别地，有“由某个流形上的辛几何确定的 Frobenius 结构，与由另一个流形上的复几何确定的 Frobenius 结构同型”。这就是经典镜像对称。

% ===== PDF p.074 (原书 p.66) =====
下面介绍一个能稍微感受到弦理论氛围的事实。设 $\overline{M}_{0,n}$ 为 $\CP^1$ 及其上 $n$ 个点组成的组的 Deligne--Mumford 模空间。此时，给出 Frobenius 结构 $(V,F_0,E,e)$，等价于考虑这样一个组 $(V,\eta,I_n)$：带有非退化双线性形式的向量空间 $(V,\eta)$，以及（带次数的）多重线性映射
\[
  I_n:V^{\otimes n}\longrightarrow H^*(\overline{M}_{0,n}),\qquad n\ge 3,
\]
其中 $I_n$ 要与嵌入 $\overline{M}_{0,n_1+1}\times\overline{M}_{0,n_2+1}\hookrightarrow\overline{M}_{0,n_1+n_2}$ 满足相容性条件。这种等价是已知的。

Frobenius 结构是 1990 年前后 Dubrovin 在研究可积系统时引入的概念。为了把二维的 $N=2$ 超对称场论与拓扑场论公理化，他基于势函数以及式 (7)\footnote{满足这个式子的、具有非退化双线性形式的结合代数称为 Frobenius 代数，Frobenius 结构这个名字正来源于此。}构建了这一结构。正如开头所述，同样的结构早已出现在斋藤恭司围绕原始形式（primitive form）的周期映射研究中。那是 1980 年前后，当时称为平坦结构。这是因为，从后一立场出发，人们注意到的是：奇点的半通用形变的底空间上存在全纯平坦的坐标\footnote{底空间是复流形，它上面的坐标关于其全纯向量场空间上的全纯 Levi-Civita 联络是平坦的。}，并据此把该结构公理化。从两种不同的动机出发到达同一个结构，这是非常有趣的事情。

\section{奇点、原始形式与平坦结构}

最终想要导出的东西已经清楚了，下面转到奇点的例子。作为最简单的例子，考虑单变量多项式 $f(x)=x^{n+1}$。$f(x)=0$ 所确定的图形，是 $n+1$ 个点汇聚于原点、重数为 $n+1$ 的一个点\footnote{为了把下面要讲的内容说准确，需要引入环层空间（ringed space）这一概念。}。这是通常意义上的光滑流形（$0$ 维，因此只是 $1$ 个点），并不是奇点。稍准确地说，一般地，多项式 $f(x_1,\dots,x_n)=0$ 所确定的图形的奇点，是指 $f=\partial f/\partial x_1=\cdots=\partial f/\partial x_n=0$ 的零点按重数计入后所成的点集。本例中 $\mathrm{d}f/\mathrm{d}x=(n+1)x^n$，所以“奇点上的函数构成的环”就是 $\mathrm{Jac}(f):=\C[x]/(x^n)$。请注意，这个环作为 $\C$ 向量空间是有限维的；并且因为 $f(x)$ 关于 $x$ 是齐次的，自然给各基底 $x^i,\,i=0,\dots,n-1$ 赋予次数 $i/(n+1)$\footnote{在本领域，习惯上把多项式的次数归一化为 $1$。}。对 $\mathrm{Jac}(f)$ 的元素 $x^i,x^j$，留数（residue）

% ===== PDF p.075 (原书 p.67) =====
由下式给出：
\begin{equation}
  (x^i,x^j)\longmapsto\frac{1}{2\pi\sqrt{-1}}\oint x^i\cdot x^j\,\frac{\mathrm{d}x}{x^n}
  \tag{13}
\end{equation}
这样就定下了 $\mathrm{Jac}(f)$ 上的非退化双线性形式。而且，把 $\circ$ 取为 $\mathrm{Jac}(f)$ 上的自然积结构时，式 (7) 也成立。也就是说，从 $f(x)$ 出发得到了 Frobenius 代数。

一般情形下，作为奇点 $f(x_1,\dots,x_n)$，我们考虑使
\begin{equation}
  \mathrm{Jac}(f):=\C[x_1,\dots,x_n]\Big/\Bigl(\frac{\partial f}{\partial x_1},\dots,\frac{\partial f}{\partial x_n}\Bigr)
  \tag{14}
\end{equation}
为有限维的那些 $f$；而作为 $\mathrm{Jac}(f)$ 上的非退化双线性形式，则考虑把留数一般化后的 Grothendieck 留数（Grothendieck's residue symbol）
\begin{equation}
  (h(x),g(x))\longmapsto \mathrm{Res}_{\C^n}
  \left[\frac{h(x)g(x)\,\mathrm{d}x_1\cdots \mathrm{d}x_n}
  {\dfrac{\partial f}{\partial x_1}\cdots\dfrac{\partial f}{\partial x_n}}\right]
  \tag{15}
\end{equation}

下面来变形奇点。这里说的变形，是指在保持某个条件\footnote{在当前情形下，保持变量个数为 $1$ 个不变，只让它变动。}的前提下连续地改变图形，最终一直动到奇点消失为止。一般的 $n+1$ 次多项式按重数计入有 $n+1$ 个根。现在考虑的例子 $f(x)=x^{n+1}$ 是 $n+1$ 次多项式，可以看作全部根都退化到 $0$ 的情形。适当地平行移动 $x$，就能让所有根之和为 $0$；因此当我们试图变形 $f(x)$ 时，它可以写成
\begin{equation}
  f_a(x)=\prod_{i=0}^{n}(x-a_i),\qquad \sum_{i=0}^{n}a_i=0
  \tag{16}
\end{equation}
$f_a=0$ 的奇点是满足 $f_a=\mathrm{d}f_a/\mathrm{d}x=0$ 的点，所以只有当某对 $i,j$ 使 $a_i-a_j=0$ 时才存在奇点，
% ===== PDF p.076 (原书 p.68) =====
这一点可以看出来。换言之，当所有根 $a_i$ 互不相同时，$f_a=0$ 没有奇点。这里，$a_i-a_j$ 所表示的 $f_a=0$ 的部分图形是 $0$ 维球面；它塌缩成 $1$ 个点，就形成了奇点。也就是说，可以看作：“曾经有一个 $0$ 维球面”这一信息发生变化，以奇异性的形式显现了出来。若从更一般的 $n$ 变量多项式出发来考虑，同样可知 $n-1$ 维球面塌缩，奇异性就出现。

设 $a_k$ 为 $a_i$ 们的 $k$ 次初等对称多项式（日文原作「基本対称式」），则可以写成
\begin{equation}
  f_a(x)=x^{n+1}+a_{n-1}x^{n-1}+\cdots+a_1x+a_0
  \tag{17}
\end{equation}
读者注意到了吗？恰好把 $\mathrm{Jac}(f)$ 作为线性空间的基底 $x^k$ 们的线性组合加到 $f(x)$ 上，就实现了很好的变形。这并非单纯的偶然。对于超曲面奇点，已知：把 $\mathrm{Jac}(f)$ 的基底乘以参数后取线性组合，再加到原来的式子 $f$ 上，就能得到称为奇点半通用形变的奇点的“好”变形。

一旦给出奇点的半通用形变 $f_a(x)$，把式 (14)(15) 中的 $f$ 换成 $f_a(x)$，就得到依赖于参数 $a_k$ 的 Frobenius 代数的族 $\mathrm{Jac}(f_a)$。特别地，通过对应 $\partial/\partial a_k\mapsto\partial f_a/\partial a_k$，有同型
\begin{equation}
  \bigoplus_{k=0}^{n-1}\C[[a_0,\dots,a_k]]\,\frac{\partial}{\partial a_k}
  \cong \mathrm{Jac}(f_a)
  \tag{18}
\end{equation}
请注意，由这个同型可知，存在与 Frobenius 代数单位元相对应的特别向量场 $\partial/\partial a_0$。然后，给 $a_k$ 赋予次数、使 $f_a(x)$ 的次数为 $1$，就定下了 Euler 向量场。

问题在于：是否存在 $\bigoplus_{k=0}^{n-1}\C[[a_0,\dots,a_k]]\partial/\partial a_k$ 的一组 $\C[[a_0,\dots,a_k]]$ 基 $(\partial/\partial t_0,\dots,\partial/\partial t_{n-1})$，使得由留数 (13)（或 (15)）给出的 $\mathrm{Jac}(f_a)$ 上的非退化双线性形式成为不依赖于参数的常数。换言之，能否在参数 $a_k$ 们的空间上取到平坦坐标系 $(t_0,\dots,t_{n-1})$。而且，只要证明存在某个势函数 $F_0(t_0,\dots,t_{n-1})$，使 $\mathrm{Jac}(f_a)$ 的积结构可由 $F_0(t)$ 在平坦坐标下的三阶微分写出，就能构成 Frobenius 结构（平坦结构）。这是最困难的地方，要经由原始形式的存在性来证明。

% ===== PDF p.077 (原书 p.69) =====
原始形式的定义及其性质无法简单叙述，很遗憾这里全部省略。

下面在简单的情形下具体写出平坦坐标与势函数。$n=1,2$ 时 $a_k$ 们本身就是平坦坐标。势函数的计算很容易，我想把它留作读者的练习。若 $n=3$，则
\[
  t_0:=a_0+\frac{a_2^2}{8},\qquad t_1:=a_1,\qquad t_2:=a_2,
\]
从而可知一般情形下 $a_k$ 们并不构成平坦坐标。% [Tutor Note: 原书式（19）第一项印作 $t_0^2t_1$；按 $\deg t_0=1$、$\deg t_1=3/4$、$\deg t_2=1/2$，该项次数为 $11/4$，与其余各项的 $5/2$ 不符（$2\cdot1+1/2=5/2$，与 $EF_0=(5/2)F_0$ 一致的是 $t_0^2t_2$）。此处照录原书。]
此时计算势函数，得到
\begin{equation}
  F_0(t_0,t_1,t_2)=\frac{1}{2}t_0^2t_1+\frac{1}{2}t_0t_1^2-\frac{1}{16}t_1^2t_2^2+\frac{1}{960}t_2^5
  \tag{19}
\end{equation}
关于 Euler 向量场
\[
  E=t_0\frac{\partial}{\partial t_0}+\frac{3}{4}t_1\frac{\partial}{\partial t_1}+\frac{1}{2}t_2\frac{\partial}{\partial t_2}
\]
有 $EF_0=(5/2)F_0$。

\section{结语}

由 $(a_0,\dots,a_{n-1})$ 给出的 $\C^n$ 成为由 $f(x)=x^{n+1}$ 给出的奇点的半通用形变的参数空间。而且，其上具有平坦结构。下面再稍微详细地看一看 $n+1$ 次式的根 $a_i$ 们所构成的空间与这个 $\C^n$ 的关系。

实际上，向量空间
\begin{equation}
  \Bigl\{(a_0,\dots,a_n)\in\C^{n+1}\ \Big|\ \sum_{i=0}^{n}a_i=0\Bigr\}
  \tag{20}
\end{equation}
同型于与 $A_n$ 型根系对应的单李代数的 Cartan 子代数 $\mathfrak{h}_{A_n}$。而置换 $n+1$ 个根的这种对称性（$n+1$ 次对称群），就是 $A_n$ 型根系相应的 Weyl 群 $W_{A_n}$。于是，通过取初等对称多项式这一操作，给出了同型
\begin{equation}
  \mathfrak{h}_{A_n}/W_{A_n}\cong\C^n
  \tag{21}
\end{equation}
也就是说，$f(x)=x^{n+1}$ 的半通用形变的空间，用根
% ===== PDF p.078 (原书 p.70) =====
系的术语就可以写成 $\mathfrak{h}_{A_n}/W_{A_n}$。而且，像前一节所说明的那样，对于非常好的一类奇点的平坦结构，也可以基于根系与 Weyl 群的不变式来求得。

% [Tutor Note: 「単純特異点」在此译作「单奇点」，指 A--D--E 型超曲面奇点（即二维 Du Val 奇点），与「单纯（simplicial）」无关；下文 Klein 正多面体的有限群与 Brieskorn 的根系描述正是指这一族奇点。]
大约 100 多年前，F. Klein 把有限群与正多面体相对应，指出了它们与二维单奇点的关联。E. Brieskorn 用李代数·根系明确地叙述了这些奇点的变形·消解。而斋藤恭司进一步扩展了这些关联，并用复解析手法与基于根系的代数手法，在奇点的半通用形变的空间上引入了平坦结构。第 11 章将介绍：有限群·根系·奇点这些乍看不同的数学领域，如何因镜像对称而紧密联系起来。

\begin{reflist}
\item[{[1]}] 江口徹，『ミラー対称性入門：ミラー対称性とは』，『数学セミナー』2006 年 10 月号．
\item[{[2]}] 深谷賢治，『ミラー対称性入門：グロモフ--ウィッテン不変量入門』，『数学セミナー』2006 年 11 月号．
\item[{[3]}] 齋藤恭司，「原始保型形式」，『数学の最先端 21 世紀への挑戦』（第 1 巻），シュプリンガー·フェアラーク東京，2002，pp.~126--159．
\item[{[4]}] B.~Dubrovin, \emph{Geometry of 2D Topological Field Theories}, Lecture Notes in Mathematics 1620, Springer, 1996, pp.~120--348.
\item[{[5]}] C.~Hertling, \emph{Frobenius manifolds and moduli spaces for singularities}, Cambridge University Press, 2002.
\item[{[6]}] Yu.~I.~Manin, \emph{Frobenius manifolds, quantum cohomology, and moduli spaces}, American Mathematical Society, 1999.
\end{reflist}
% 本节号体例至此结束：恢复全书默认的罗马数字，供后续章节使用。
\renewcommand{\thesection}{\Roman{section}}
% 原书此处有作者读音标记［たかはし あっし］，按 STYLE_GUIDE 省略，记入 reports/ch05.md。

